\newpage

\section{QR-Decomposition For Numerical Mathematics}

We want to find \(A = QR\), where \(R\) is an upper-triangular matrix and \(Q\) just a matrix. 

\subsection{Givens Matrices}

This algorithm uses rotations for the elimination of the elements underneath of the main diagonal. We use 
matrices of the form:

\[
        Q_{i,j} = 
        \begin{bmatrix}
            1 &        &         &        &         & & &\\
              & \ddots &         &        &         & & &\\
              &        & c       &        &   s    & & &\\
              &        &         & \ddots &         & & &\\
              &        & -s      &        &   c     & & &\\
              &        &         &        &         & \ddots & &\\
              &        &         &        &         &  & & 1
        \end{bmatrix} 
\]

where given an element \(a_{k,l}\) to eliminate then 

\begin{align*}
    q_{k - 1, k - 1} &= c \\ 
    q_{k - 1, k} &= s, \\
    q_{k, k - 1} &= -s,\\ 
    q_{k, k} &= c
\end{align*}

and 

\[
    s = \frac{a_{k, l}}{\sqrt{a_{k - 1, l}^2 + a_{k,l}^2}}, \quad c = \frac{a_{k - 1, l}}{\sqrt{a_{k - 1,l}^2 + a_{k,l}^2}}
\]

Which when chained will give us 

\[
    Q_{n-1}\cdot \dots \cdot Q_1 A = R
\]

By the orthonormality of the \(Q_i\) we have 

\[
    A = Q_1^T \cdot \dots \cdot Q_{n-1}^T R = Q R
\]

\subsubsection{Givens Algorithm}

We have our system \(Ax = b\) in \(\R^{n \times n}\) and are going to store the \(Q\) and \(R\) in the same matrix. 
We start with \(Q = I\) and \(R = A\).

\begin{enumerate}

    \item For each element \(a_{i, j}\) underneath the main diagonal we compute: 

    \[
        s = \frac{a_{i,j}}{\sqrt{a_{i,j}^2 + a_{k,j}^2}}, \quad c = \frac{a_{k, j}}{\sqrt{a_{i,j}^2 + a_{k,j}^2}}
    \]

    where \(k\) is the \textbf{index of the first non-zero entry in the current column above} \(a_{i, j}\).

    \item Create the matrix \(Q_{i,j}\).

    \item Multiply \(Q_{i,j}\) with \(A\) and store at the position of the generated zero the value \(s\) and at the position of \(c\) the value \(c\).

    \item Continue with the next element in the column until all elements are eliminated. 
    Note that we applying the multiplication for the new \(A\), we handle the places where we have stored the \(\rho_{i - 1, j - 1}\) as if they were zero, 
    because they are not used for the next steps.
    
    \item Then, continue with the next column in the new sub-matrix and so on until we have an upper triangular matrix.

    \item Now, compute \(y = Q^T b = Q_{n-1} \cdot \dots \cdot Q_1 b\). One matrix a a time. We use the stored coefficient \(\rho_{i,j} = s\) to then compute 
    \(s = \sqrt{1 - s^2}\) for the multiplication.

    \item Finally, solve \(Rx = Q_{n-1} \cdot \dots \cdot Q_1 Ax= y\) for \(x\) using back substitution.

\end{enumerate}

It is important to note that this algorithm is four times more expensive that the \(LU\)-Decomposition, but it is more stable.

\textbf{Example:}

Find the \(QR\)-Decomposition of \(A\) and solve \(Ax = b\). 

\[
    A = 
    \begin{bmatrix}
        12 & -40 & 46 \\ 
        -9 & 30 & 153 \\ 
        20 & 100 & 135
    \end{bmatrix}
    , 
    b = 
    \begin{bmatrix}
       -16 \\ 
       12 \\ 
       140
    \end{bmatrix}
\]


We start by eliminating \(a_{2,1}\): 

\[
    c = \frac{12}{\sqrt{12^2 + (-9)^2}} = \frac{4}{5}, \ 
    s = \frac{-9}{\sqrt{12^2 + (-9)^2}} = - \frac{3}{5}
\]

Therefore 

\[
    Q_{2,1} =
    \begin{bmatrix}
        \frac{4}{5} & - \frac{3}{5} & 0 \\ 
        \frac{3}{5} & \frac{4}{5} & 0 \\ 
        0 & 0 & 1
    \end{bmatrix}
\]

Then

\[
    Q_{2,1}A =
    \begin{bmatrix}
        15 & -50 & -55 \\ 
        0 & 0 & 150 \\ 
        20 & 100 & 135
    \end{bmatrix}
\]

We store \(\rho_{2,1} = s = - \frac{3}{5}\) at the generated 0 at \(a_{2,1}\).

\[
    \begin{bmatrix}
        15 & -50 & -55 \\ 
        - \frac{3}{5} & 0 & 150 \\ 
        20 & 100 & 135
    \end{bmatrix}
\]

Now, we eliminate the element \(a_{3, 1}\):

\[
    c = \frac{15}{\sqrt{15^2 + 20^2}} = \frac{3}{5}, \ 
    s = \frac{20}{\sqrt{15^2 + 20^2}} = \frac{4}{5}
\]

Now we get 

\[
    Q_{3,1} = 
    \begin{bmatrix}
        \frac{3}{5} & 0 & \frac{4}{5} \\
        0 & 1 & 0 \\
        -\frac{4}{5} & 0 & \frac{3}{5}
    \end{bmatrix}
\]

Applying the multiplication to our previously computed matrix \(Q_{2,1}A\)

\[
    Q_{3,1}Q_{2,1}A = 
    \begin{bmatrix}
        25 & 50 & 75 \\ 
        0 & 0 & 150 \\ 
        0 & 100 & 125
    \end{bmatrix}
\]

We store \(\rho_{3,1} = s = \frac{4}{5}\) at \(a_{3,1}\). 

\[
    \begin{bmatrix}
        25 & 50 & 75 \\ 
        -\frac{3}{5} & 0 & 150 \\ 
        \frac{4}{5} & 100 & 125
    \end{bmatrix}
\]

Continuing with \(a_{3, 2}\):

\[
    c = \frac{0}{\sqrt{0^2 + 100^2}} = 0, \ 
    s = \frac{100}{\sqrt{0^2 + 100^2}} = 1 
\]

\[
    Q_{3,2} = 
    \begin{bmatrix}
        1 & 0 & 0 \\
        0 & 0 & 1 \\
        0 & -1 & 0
    \end{bmatrix}
\]

Multiply 

\[
    Q_{3,2}Q_{3,1}Q_{2,1}A =
    \begin{bmatrix}
        25 & 50 & 75 \\ 
        0 & 100 & 125 \\ 
        0 & 0 & -150
    \end{bmatrix}
\]

Now, we store \(\rho_{3,1} = s = 1\) at \(a_{3,2}\)

\[
    \begin{bmatrix}
        25 & 50 & 75 \\ 
        -\frac{3}{5} & 100 & 125 \\ 
        \frac{4}{5} & 1 & -150
    \end{bmatrix}
\]

Now we solve \(Qy = b \iff Q^T b = y\) for \(Rx = y\) using 

\[
    Q = (Q_{3,2}Q_{3,1}Q_{2,1})^T
    = 
    \frac{1}{25}
    \begin{bmatrix}
        12 & -16 & -15 \\ 
        -9 & 12 & -20 \\ 
        20 & 15 & 0
    \end{bmatrix}
\]

and 

\[
    R = Q_{3,2}Q_{3,1}Q_{2,1}A = 25
    \begin{bmatrix}
       1 & 2 & 3 \\ 
       0 & 4 & 5 \\ 
       0 & 0 & -6 
    \end{bmatrix}
\]

For each of the matrices in \(Q^Tb\) we reconstruct the the \(Q_{i,j}\) matrix and perform the 
multiplication 

Using \(\rho_{2,1} = -\frac{3}{5}, c = \sqrt{1 - s^2} = \frac{4}{5}\):

\[
    Q_{2,1}b = 
    \begin{bmatrix}
        \frac{4}{5} & - \frac{3}{5} & 0 \\ 
        \frac{3}{5} & \frac{4}{5} & 0 \\ 
        0 & 0 & 1
    \end{bmatrix}
    \begin{bmatrix}
        -16 \\ 
        12 \\ 
        140
    \end{bmatrix}
    = 
    \begin{bmatrix}
        -20 \\ 
        0 \\ 
        140
    \end{bmatrix}
\]

Then repeat for the others, and apply the multiply the new \(b\) to the reconstructed matrix and we get 

\[
    y = 
    \begin{bmatrix}
        100 \\ 
        100 \\ 
        0
    \end{bmatrix}
\]

Finally, for \(Rx = y\) we have 

\[
    25 
    \begin{bmatrix}
       1 & 2 & 3 \\ 
       0 & 4 & 5 \\ 
       0 & 0 & -6 
    \end{bmatrix}
    x 
    = 
    y
\]

which gives us \(x = [2, 1, 0]^T\).

\subsection{Householder Algorithm}

The \emph{Householder algorithm} similarly to Givens algorithm will given us a decomposition of a matrix \(A\) 
such that \(A = QR\) with \(Q\) being orthonormal and \(R\) being upper triangular. The difference relies on the method of the transformations 
which instead of being rotations, they are reflections.

\subsubsection{Householder Transformation}

Given a vector \(v \in \R^n\) we denote 

\[
    H_v = I - 2 \frac{vv^T}{v^Tv}
\]

as the \emph{Householder transformation} for \(v\). Equivalently 

\[
    H_v = I - 2 \frac{vv^T}{\metric{v}_2^2}.
\]

If \(v\) is normalized, then \(H_v = I - 2 vv^T\), which is a reflection across the hyperplane orthogonal to \(v\).

\textbf{Example:}

Given \(v = [1, 2]^T\) we have

\[
    H_v = 
    \begin{bmatrix}
        1 & 0 \\ 
        0 & 1
    \end{bmatrix}
    - 2
    \frac{
        \begin{bmatrix}
            1 \\ 2
        \end{bmatrix}
        \begin{bmatrix}
            1 & 2
        \end{bmatrix}
    }{\begin{bmatrix}
        1 & 2
    \end{bmatrix}    
    \begin{bmatrix}
        1 \\ 2
    \end{bmatrix}}
    =
    \begin{bmatrix}
        1 & 0 \\ 
        0 & 1
    \end{bmatrix}
    - 2 \frac{1}{5}
    \begin{bmatrix}
        1 & 2 \\ 
        2 & 4
    \end{bmatrix}
    =
    \frac{1}{5}
    \begin{bmatrix}
        -3 & -4 \\ 
        -4 & 3
    \end{bmatrix}
\]

\subsection{Effect of a Householder Transformation}

Given \(x \in \R^n\) when multiplying \(H_k\), which is the \(k\)-th Householder transformation,
with \(x\) we get

\[
    H_k x =
    H_k
    \begin{bmatrix}
        x_1 \\ 
        x_2 \\ 
        \vdots \\
        x_{k - 1} \\ 
        x_k \\ 
        x_{k + 1} \\ 
        \vdots \\
        x_n
    \end{bmatrix}
    = 
    \begin{bmatrix}
        x_1 \\ 
        x_2 \\ 
        \vdots \\
        x_{k - 1} \\     
        ? \\ 
        0 \\ 
        \vdots \\
        0
    \end{bmatrix}
\]

The vector to be used for the transformation defined as 

\[
    v = x \pm \metric{x}_2 e_k
\]

The sign is not really important, but it is usually chosen to avoid cancellation errors. 

The \(k\)-th Householder transformation is then defined as \(\tilde{x}\) where all elements before the \(k\)-th element are zero, i.e. if we have to eliminate 
all elements below the first element, then we have \(k = 1\) an \(\tilde{x} = x\), but if we have to eliminate all elements below the second element, then we have \(k = 2\) and \(\tilde{x} = [0, x_2, x_3, \dots, x_n]^T\) and so on.
Note that the resulting transformation will be still be used with \(x\).

\subsubsection*{Householder Algorithm}

Given \(A \in \R^{n \times n}\) we start with \(Q = I\) and \(R = A\).
Similar to the \(LU\)-Decomposition we will start with the first column and then focus on the sub-matrix and so on.

\begin{enumerate}

    \item Eliminate the elements below the pivot element \(a_{i,i}\) by applying the Householder transformation \(H_i\) to the \(i\)-th column of the current matrix.

    \item Multiply \(H_i\) with the current matrix and store the resulting matrix which will be used for the next step.

    \item Continue with the next column and so on until we have an upper triangular matrix. The result should be \(H_k \cdot H_{k - 1} \cdot \dots \cdot H_1 \cdot A = R\) and \(Q = H_1 \cdot H_2 \cdot \dots \cdot H_k\).
 
\end{enumerate}

Note that the Householder matrices are self-inverse, i.e. \(H_i^{-1} = H_i\) and therefore \(A = Q R\).

\textbf{Example}

Given is  

\[
    A = 
    \begin{bmatrix}
        1 & 2 \\ 
        3 & 4
    \end{bmatrix}
\]

We start with the first column and compute the Householder transformation for the first column. We have

\[
    v = 
    \begin{bmatrix}
        1 \\ 
        3
    \end{bmatrix}
    + \sqrt{10} 
    \begin{bmatrix}
        1 \\ 
        0
    \end{bmatrix}
    =
    \begin{bmatrix}
        1 + \sqrt{10} \\ 
        3
    \end{bmatrix}
\]

Then, we compute the Householder transformation

\[
    H_1 = I - 2 \frac{vv^T}{v^Tv} = 
    \begin{bmatrix}
        -\frac{1}{5} & -\frac{3}{5} \\ 
        -\frac{3}{5} & \frac{1}{5}
    \end{bmatrix}
\]

And we apply it to \(A\)

\[
    H_1 A = 
    \begin{bmatrix}
        -\frac{1}{5} & -\frac{3}{5} \\ 
        -\frac{3}{5} & \frac{1}{5}
    \end{bmatrix}
    \begin{bmatrix}
        1 & 2 \\ 
        3 & 4
    \end{bmatrix}
    =
    \begin{bmatrix}
        -\sqrt{10} & -\frac{4}{5}\sqrt{10} \\ 
        0 & \frac{2}{5}\sqrt{10}
    \end{bmatrix}
    = R
\]

Now, we have an upper triangular matrix and we can compute \(Q\) as \(H_1\) and \(R\) as \(H_1 A\), and we have \(A = QR\).

\section{Iterative Methods}

An \emph{iterative method} 

\[
    x^{(k + 1)} = \phi_k(x^{(k)}, x^{(k - 1)}, \dots, x^{(0)})
\]

is convergent if \(x^{(k)} \to x\) for all  \(x^{(0)} \in \R^n\). It also has the \emph{fixed point property} 
\(x = \phi_(x) \iff Ax = b\).

The \emph{error} \(e^{(k)}\) and the \emph{residuum} \(r^{(k)}\) are defined as

\[
    e^{(k)} = x - x^{(k)}, \quad r^{(k)} = b - Ax^{(k)}
\]

If for the computation of \(x^{(k + 1)}\) only the direct predecessor \(x^{(k)}\) uses, then

\[
    x^{(k + 1)} = \phi(x^{(k)})
\]

the process is considered with \emph{one step} \(\phi_k : \R^n \to \R^n\), else is considered with \emph{multi step} 
\(\phi_k : (\R^n)^{k + 1} \to \R^n\).

If \(\phi_k = \phi\) for all \(k\) then the process is called \emph{stationary}, else it is called \emph{non-stationary}.

\subsection{Single Step Stationary Iteration Method}

A \emph{single step stationary iteration method} is given by 

\[
    x^{(k + 1)} = \phi(x^{(k)}) = Bx + c,
\]

where \(B\) is the \emph{iteration matrix} and \(c\) is the \emph{iteration vector}.

\subsection{Error of Single Step Stationary Iteration Method}

If the single step stationary iteration method has the fixed point property, then for 
\(e^{(k)} = x - x^{(k)}, r^{(k)} = b - Ax^{(k)}\) we have

\[
    Ae^{(k)} = r^{(k)}
\]

\textbf{Proof:}

\[
    Ae^{(k)} = A(x - x^{(k)}) = Ax - Ax^{(k)} = b - Ax^{(k)} = r^{(k)}
\]

\QED

\subsection{Convergence}

Analyzing our regular linear system \(Ax = b\) with the single step iteration 

\[
    x^{(k + 1)} = \phi(x^{(k)}) = Bx^{(k)} + c
\]

with the fixed point property \(x = \phi(x) = Bx + c \iff Ax = b\) we get

For the error \(e^{(k)} = x^{(k)} - x\) we have 

\[
    e^{(k + 1)} = x^{(k + 1)} - x = Bx^{(k)} + c - (Bx + c) = Be^{(k)}
\]

Therefore, 

\[
    e^{(k)} = B^k e^{(0)}
\]

Then 

\[
    e^{(k)} \to 0 \quad \text{as} \quad k \to \infty
\]

for an arbitrary \(e^{(0)}\) if and only if \(B^k \to 0 \quad \text{as} \quad k \to \infty\).

\subsection{The Spectral Radius}

The \emph{spectral radius} of a matrix \(B\) is defined as
 
\[
    \rho(B) = \max\{|\lambda| \mid \lambda \in \sigma(B)\},
\]

where \(\sigma(B)\) is the \emph{spectrum} of \(B\), i.e. the set of all eigenvalues of \(B\).

\subsubsection*{Properties of the Spectral Radius}

\begin{itemize}

    \item \(\rho(B) \le \metric{B}_I\) for all induced metrics.

    \item \(\forall \varepsilon > 0\) there exists an induced metric \(\metric{\cdot}_I\) from a vector metric \(\metric{\cdot}_V\) such that:
    \[
        \metric{B}_I = \sup_{w \ne 0} \frac{\metric{Bw}_w}{\metric{w}_V}
    \]
    with \(\metric{B}_I \le \rho(B) + \varepsilon\).

\end{itemize}

\subsection{Gelfand}

For any matrix metric \(\metric{\cdot}\) the following holds:

\[
    \rho(B) = \lim_{k \to \infty} \sqrt[k]{\metric{B^k}}.
\]

Also for an induced matrix metric \(\metric{\cdot}_I\) \(\metric{B}_I^{\frac{1}{k}}\) converges towards \(\rho(B)\).
\subsection{Splitting Method For Linear Systems}

Given a regular matrix \(A \in \R^{n \times n}\) in a system \(Ax = b\), we can split \(A\) into 

\[
    A = M - N
\]

where \(M\) is regular and \(N\) is arbitrary. Then, we can write

\[
    Mx = Nx + b
\]

and we can define the iteration 

\[
    x^{(k + 1)} = M^{-1}Nx^{(k)} + M^{-1}b,
\]

and we set \(x^{(0)} = \vec{0} \in \R^n\) as the initial guess.

\textbf{Example:}

Given the system 

\[
    \begin{bmatrix}
        4 & -1 & 0 & 0 \\ 
        -1 & 4 & -1 & 0 \\ 
        0 & -1 & 4 & -1 \\ 
        0 & 0 & -1 & 3
    \end{bmatrix}
    x = 
    \begin{bmatrix}
        15 \\ 
        10 \\ 
        10 \\ 
        10
    \end{bmatrix}
\]

Then we can split \(A\) into

\[
    M = 
    \begin{bmatrix}
        4 & 0 & 0 & 0 \\ 
        0 & 4 & 0 & 0 \\ 
        0 & 0 & 4 & 0 \\ 
        0 & 0 & 0 & 3
    \end{bmatrix}
    , \quad 
    N = 
    \begin{bmatrix}
        0 & -1 & 0 & 0 \\ 
        -1 & 0 & -1 & 0 \\ 
        0 & -1 & 0 & -1 \\ 
        0 & 0 & -1 & 0
    \end{bmatrix}
\]

Which results in the iteration

\[
    x^{(k + 1)} = 
    \begin{bmatrix}
        \frac{1}{4} & 0 & 0 & 0 \\ 
        0 & \frac{1}{4} & 0 & 0 \\ 
        0 & 0 & \frac{1}{4} & 0 \\
        0 & 0 & 0 & \frac{1}{3} 
    \end{bmatrix}
    \begin{bmatrix}
        0 & -1 & 0 & 0 \\ 
        -1 & 0 & -1 & 0 \\ 
        0 & -1 & 0 & -1 \\ 
        0 & 0 & -1 & 0
    \end{bmatrix}
    \begin{bmatrix}
       0 \\ 
       0 \\ 
        0 \\
        0 
    \end{bmatrix}
    + 
    \begin{bmatrix}
       \frac{1}{4} & 0 & 0 & 0 \\ 
        0 & \frac{1}{4} & 0 & 0 \\ 
        0 & 0 & \frac{1}{4} & 0 \\
        0 & 0 & 0 & \frac{1}{3} 
    \end{bmatrix}
    \begin{bmatrix}
        15 \\ 
        10 \\ 
        10 \\ 
        10
    \end{bmatrix}
    = 
    \begin{bmatrix}
        3.75 \\ 
        2.5 \\ 
        2.5 \\ 
        3.333\ldots
    \end{bmatrix}
\]

We can then continue the process using our resulting vector as the new \(x^{(k)}\) and we will get closer to the solution of the system.

\subsubsection*{Triple Splitting}

We can also define the splitting 

\[
    A = D - L - U
\]

where \(D\) is the diagonal of \(A\), \(L\) is the lower triangular part of \(A\) and \(U\) is the upper triangular part of \(A\). 
We then derive from that 

\begin{align*}
    Ax = b &\iff (D - L - U)x = b \\ 
           &\iff Dx = (L + U)x + b
\end{align*}

The iteration is then given by

\[
    x^{(k + 1)} = D^{-1}(L + U)x^{(k)} + D^{-1}b
\]

\subsection{The Residuum As Stopping Criterion}

The \emph{residuum} is defined as

\[
    r^{(k)} = b - Ax^{(k)}
\]

When \(\metric{r^{(k)}} \le \varepsilon \metric{b}\) for a given \(\varepsilon > 0\) then we can stop the iteration because we are close enough to the solution.

\subsection{Convergence For Semi Positive Definite Matrices}

Given semi positive definite matrices \(A, M\); if \(2M - A\) is also semi positive definite, then 

\[
    \phi(x) = Bx + d,\ B = I - M^{-1}A, \ d = M^{-1}b
\]

converges and also \(\metric{B}_A = \metric{B}_M = \rho(B) < 1\).

Also, the splitting \(M + M^T - A\) is also spd. and \(\rho(B) \le \metric{B}_A < 1\).

\subsection{Strict Diagonal-dominant}

\(A\) is strictly diagonal-dominant if 

\[
    |a_{i,i}| > \sum_{j \ne i} |a_{i,j}|, \ \forall i = 1, \dots, n
\]

\subsection{The Jacobi Method}

The Jacobi method is a single step stationary iteration method with the splitting 
\(A = D - (L + U)\) and the iteration:

\[
    x^{(k + 1)} = D^{-1}(L + U)x^{(k)} + D^{-1}b
\]

\textbf{Example:}

Given the system 

\[
    \begin{bmatrix}
        4 & -1 & 0 & 0 \\ 
        -1 & 4 & -1 & 0 \\ 
        0 & -1 & 4 & -1 \\ 
        0 & 0 & -1 & 3
    \end{bmatrix}
    b = 
    \begin{bmatrix}
        15 \\ 
        10 \\ 
        10 \\ 
        10
    \end{bmatrix}
\]

We can apply the Jacobi method with the splitting \(A = D - (L + U)\) where

\[
    D = 
    \begin{bmatrix}
        4 & 0 & 0 & 0 \\ 
        0 & 4 & 0 & 0 \\ 
        0 & 0 & 4 & 0 \\ 
        0 & 0 & 0 & 3
    \end{bmatrix}
    , \quad 
    L + U = 
    \begin{bmatrix}
        0 & -1 & 0 & 0 \\ 
        -1 & 0 & -1 & 0 \\ 
        0 & -1 & 0 & -1 \\ 
        0 & 0 & -1 & 0
    \end{bmatrix}
\]

Then, we can compute the iteration

\[
    x^{(k + 1)} = 
    \begin{bmatrix}
        \frac{1}{4} & 0 & 0 & 0 \\ 
        0 & \frac{1}{4} & 0 & 0 \\ 
        0 & 0 & \frac{1}{4} & 0 \\
        0 & 0 & 0 & \frac{1}{3} 
    \end{bmatrix}
    \begin{bmatrix}
        0 & -1 & 0 & 0 \\ 
        -1 & 0 & -1 & 0 \\ 
        0 & -1 & 0 & -1 \\ 
        0 & 0 & -1 & 0
    \end{bmatrix}
    \begin{bmatrix}       
        0 \\ 
        0 \\ 
        0 \\
        0 
    \end{bmatrix}
    +
    \begin{bmatrix}
         \frac{1}{4} & 0 & 0 & 0 \\ 
          0 & \frac{1}{4} & 0 & 0 \\ 
          0 & 0 & \frac{1}{4} & 0 \\    
          0 & 0 & 0 & \frac{1}{3} 
    \end{bmatrix}
    \begin{bmatrix}
        15 \\ 
        10 \\ 
        10 \\ 
        10
    \end{bmatrix}
\]

Giving us

\[
    x^{(1)} = 
    \begin{bmatrix}
        3.75 \\ 
        2.5 \\ 
        2.5 \\ 
        3.333\ldots
    \end{bmatrix}
\]

\subsection{Gauss-Seidel Method}

The Gauss-Seidel method is a single step stationary iteration method with the splitting
\(A = (D - L) - U\) and the iteration

\[
    x^{(k + 1)} = (D - L)^{-1}Ux^{(k)} + (D - L)^{-1}b
\]

\subsection{After Iteration}

This process refers to using an iterative method after using a mathematical algorithm for solving a certain problem to 
mitigate rounding errors. In our case we will present such a process with the \(LU\)-Decomposition.

While performing the \(LU\)-Decomposition we will get \(\tilde{L}\tilde{U} = M\) which are the \(L,U\) matrices, but with rounding errors.
We use the stationary linear iteration to improve the solution \(\tilde{x}\) for which \(A\tilde{x} = b\).

\begin{enumerate}

    \item Perform the \(LU\)-Decomposition to get the triangular matrices and the approximate solution with errors.

    \item Set \(x_0 = \tilde{x}\) 

    \item Repeat the following steps as long \(\frac{\metric{p_k}}{\metric{x_{k + 1}}} < \varepsilon\): 
    
    \begin{itemize}
        \item \(r_k = b - Ax_k \).
        \item \(p_k = M^{-1}r_K\) in other words \(\tilde{L}\tilde{U}p_k = r_k\). Also \begin{verbatim}p_k = np.linalg.solve(A, r_k)\end{verbatim}
        \item \(x_{k + 1} = x_k + p_k\).
    \end{itemize}
\end{enumerate}

\subsection{Gradient Descent}

Given our stationary iteration 

\[
    x^{(k + 1)} = x^{(k)} + M^{-1}r^{(k)}
\]

We add the parameter \(\alpha^{(k)}\)

\[
    x^{(k + 1)} = x^{(k)} + \alpha^{(k)}M^{-1}r^{(k)}
\]

With the goal in each step to find an \(\alpha^{(k)}\) which gives us the "best" \(x^{(k + 1)}\).
This means 

\begin{align*} 
    e^{(k + 1)} &= x - x^{(k + 1)} \\ 
    &= x - x^{(k)} + \alpha^{(k)}M^{-1}r^{(k)} \\ 
    &= e^{(k)} - \alpha^{(k)}p^{(k)}
\end{align*}

Geometrically, \(e^{(k + 1)}\) is on a straight line determined by \(e^{(k)}\) in the direction \(p^{(k)}\). This means 
that \(\metric{e^{(k + 1)}}\) will be minimal when \(p_{(k)}\) is orthogonal to \(e^{(k + 1)}\)

\begin{align*}
    \inner{e^{(k + 1)}}{p^{(k)}} &= \inner{p^{(k)}}{e^{(k)} - \alpha^{(k)}p^{(k)}} \\ 
    &= \inner{p^{(k)}}{e^{(k)}} - \alpha^{(k)} \inner{p^{(k)}}{p^{(k)}} \\ 
    \alpha^{(k)} &= \frac{\inner{p^{(k)}}{e^{(k)}}}{\inner{p^{(k)}}{p^{(k)}}}
\end{align*}

Hence, we have constructed the following iterative algorithm 

\begin{align*}
   r^{(k)} &= b - Ax^{(k)}  \\
   p^{(k)} &= M^{-1}r^{(k)} \\ 
   \alpha^{(k)} &= \frac{\inner{p^{(k)}}{e^{(k)}}}{\inner{p^{(k)}}{p^{(k)}}} \\ 
   x^{(k + 1)} &= x^{k} + \alpha^{(k)}p^{(k)}
\end{align*}

The problem is that technically we do not know the exact solution \(x\), but this can be solved using a 
suited inner product with metric 

\[
    \metric{x} = \metric{x}_A = \sqrt{x^T A x}
\]

Giving us 

\[
   \alpha^{(k)} =  \frac{\inner{p^{(k)}}{e^{(k)}}_A}{\inner{p^{(k)}}{p^{(k)}}_A}
   = \frac{\inner{p^{(k)}}{Ae^{(k)}}}{\inner{p^{(k)}}{Ap^{(k)}}}
   = \frac{\inner{p^{(k)}}{r^{(k)}}}{\inner{p^{(k)}}{Ap^{(k)}}}
\]

\section{Iterative Methods For Eigenvalue Problems}

\subsection{The Vector Iteration}

Given a normal matrix \(A \in \C^{n \times n}\), which is unitarily diagonalizable, i.e. there exists an orthonormal basis of 
eigenvectors, \(u_1, \dots, u_n\).

Each vector \(x \in \C^n\) can be written as \(x = \sum_i x_i u_i\) and 

\[
    Ax = \sum_i x_i \lambda_i u_i, \quad A^k x = \sum_i x_i \lambda_i^k u_i
\]

If \(|\lambda_1| > |\lambda_2| \ge \dots \ge |\lambda_n|\) then the dominant eigenvalue is \(\lambda_1\). This 
means that by applying \(A\) to \(x\) multiple times we get closer to the dominant eigenvector \(u_1\) and the dominant eigenvalue \(\lambda_1\).

\subsubsection*{Algorithm}

\begin{enumerate}
    
    \item Choose a \(x^{(0)}\) such that \(\metric{x^{(0)}}_2 = 1\) and \(\inner{x^{(0)}}{u_1} \neq 0\).

    \item Repeat the following steps until convergence: 
    
    \begin{align*}
        y^{(k + 1)} &= Ax^{(k)} \\
        x^{(k + 1)} &= \frac{y^{(k + 1)}}{\metric{y^{(k + 1)}}_2} \\ 
        \mu^{(k + 1)} &= \inner{x^{(k + 1)}}{Ax^{(k + 1)}} = \frac{\inner{x^{(k + 1)}}{y^{(k + 1)}}}{\metric{x^{(k)}}_2^2}
    \end{align*}

\end{enumerate}

\subsection{The Jacobi Method}

The \emph{Jacobi method} is an iterative method for computing the eigenvalues of a symmetric matrix \(A\) by applying a sequence 
of orthogonal transformations to \(A\) to make it diagonal. 

The eigenvalues of \(A\) are then given by the diagonal elements of the resulting matrix.

The derivation will come from applying orthonormal matrices repeatedly to to our original matrix. 

\[
    A^{(0)} = A, \quad A^{(k)} = Q_k^T A^{(k - 1)} Q_k
\]

Using our Givens-Rotations 

\[
    Q_k = 
    \begin{bmatrix}
        1 &        &    &         &        &        &         &   \\
          & \ddots &    &         &        &        &         &   \\
          &        & c  &         &        &   s    &         &   \\
          &        &    &   1     &        &        &         &   \\
          &        &    &         & \ddots &        &         &   \\
          &        & -s &         &        &   c    &         &   \\
          &        &    &         &        &   1    &         &   \\
          &        &    &         &        &        & \ddots  &   \\
          &        &    &         &        &        &         & 1 \\
    \end{bmatrix}
\]


where 


\[
   c = \cos(\alpha), \quad s = \cos(\alpha), \alpha = 
   \begin{cases}
    \frac{\pi}{4}  & a_{i_0, i_0} = a_{j_0, j_0} \\ 
    \frac{1}{2}\arctan\left(\frac{2a_{i_0, j_0}}{a_{j_0, j_0} - a_{i_0, i_0}}\right)
   \end{cases} 
\]

Note the following coordinates when reading the matrix from top to bottom, left to right:

\begin{align*}
    c & \quad \text{at} \quad i_0, i_0 \\ 
    s & \quad \text{at} \quad i_0, j_0 \\
    c & \quad \text{at} \quad j_0, j_0 \\
    -s & \quad \text{at} \quad j_0, i_0
\end{align*}

And we use the iteration 

\[
    Q^{(0)} = I, \quad Q^{(k)} = Q^{(k - 1)}Q_k
\]

\subsubsection*{Stop Criterion}

We define the square sum of the diagonal elements of the matrices \(A^{(k)}\) as 

\[
    N(A^{(k)}) = \sum_{i \ne j} \left(a_{i,j}^{(k)}\right)^2.
\]

It holds that 

\[
    N^(A^{(k + 1)}) \le \alpha N(A^{(k)}), \quad \alpha \in [0, 1)
\]

converges \emph{strictly monotonously decreasing} towards zero. Thus, this is our convergence criteria 
given an upper bound.

\subsection{The QR Method}

We take a matrix \(A \in \R^{n \times n}\) with \(n\) start vectors whose dimension is also \(n\).

Choose \(X^{(0)} = [x_1^{(0)},\dots, x_n^{(0)}]\) where the vectors form an orthonormal basis

For \(k = 0, 1, \dots\) repeat the following steps:

\begin{align*}
    Y^{(k + 1)} &= AX^{(k)} \\ 
    \tilde{Q^{(k + 1)}} R^{(k + 1)} &= Y^{(k + 1)} \\ 
    X^{(k + 1)} &= \tilde{Q^{(k + 1)}} \\
    T^{(k + 1)} &= X^{(k)^T} Y^{(k + 1)} = X^{(k)^T} A X^{(k)}
\end{align*}

Then for 

\begin{align*} 
    T^{(k)} &= X^{(k - 1)^T} A X^{(k - 1)} \\ 
    &= X^{(k - 1)^T} Y^{(k)} \\ 
    &= X^{(k - 1)^T} \tilde{Q^{(k)}} R^{(k)}\\ 
    &= (X^{(k - 1)^T} X^{(k)}) R^{(k)} \\ 
    &= Q^{(k)} R^{(k)}
\end{align*}

We can also use \(Q^{(k)} = X^{(k - 1)^T} X^{(k)}\). And because the vector of \(X^{(k)}\) are orthonormal we also have that 

\begin{align*} 
    T^{(k + 1)} &= X^{(k)^T} A X^{(k)}\\
    &=  X^{(k)^T} A X^{(k - 1)} X^{(k - 1)^T} X^{(k)}\\
    &=  X^{(k)^T} Y^{(K)} X^{(k - 1)^T} X^{(k)}\\
    &=  X^{(k)^T} \tilde{Q^{(k)}} R^{(k)} X^{(k - 1)^T} X^{(k)}\\
    &=  X^{(k)^T} X^{(k)} R^{(k)} X^{(k - 1)^T} X^{(k)}\\
    &=  R^{(k)} (X^{(k - 1)^T} X^{(k)})\\
    &=  R^{(k)} Q^{(k)} 
\end{align*}

This means that our iteration can be simplified to 

\[
    A^{(k)} = Q^{(k)}R^{(k)}, \quad A^{(k + 1)} =  R^{(k)} Q^{(k)} 
\]

Because \(A\) is in the real Schur-form \(A = Q T Q^T\), under certain circumstances we get: 

\[
    A^{(k)} \rightarrow^{k \to \infty} T, \quad \hat{Q}^{(k)} \rightarrow^{k \to \infty} T
\]

with \(\hat{Q}^{(k)} = Q^{(0)} \cdot \dots \cdot Q^{(k - 1)}\). 

